\chapter{Value-Based Deep Reinforcement Learning - DQN}
\label{chap:dqn}

Chapter~\ref{chap:reinforcementlearningfoundations} identified two paths to an optimal policy. This chapter develops the first of them, the value-based route which estimates the optimal action-value function $q_*$ and derives a policy greedily from the estimate. In finite state spaces, the tabular Q-learning algorithm \citep{Watkins1989LearningFD} realizes this idea. Pong's observation space, however, is far too large for a tabular representation to be viable. The value function must therefore be approximated, and the resulting algorithm does not inherit tabular Q-learning's convergence guarantees \citep[Ch.~11]{suttonReinforcementLearningIntroduction2020}. Section~\ref{sec:fromqlearningtodqn} traces the path from tabular Q-learning to its deep, function-approximating counterpart and states the resulting optimization objective. Section~\ref{sec:stabilization} examines why a direct implementation of this objective is unstable and presents experience replay and target networks as the two mechanisms through which \acrshort{dqn} \citep{mnihHumanlevelControlDeep2015} restores stable learning.

\section{From Q-Learning to Deep Q-Networks}
\label{sec:fromqlearningtodqn}
Tabular Q-learning is an off-policy \acrshort{td} control algorithm defined by:
\begin{equation}\label{eq:q_learning}
	Q\left(S_t, A_t\right) \leftarrow Q\left(S_t, A_t\right) + \alpha \left[R_{t+1} + \gamma \max_a Q\left(S_{t+1}, a\right) - Q\left(S_t, A_t\right)\right].
\end{equation}
$Q$ is the learned (updated) action-value function that directly approximates the Bellman optimal action-value function $q_*$. The bracketed quantity $R_{t+1} + \gamma \max_a Q\left(S_{t+1}, a\right) - Q\left(S_t, A_t\right)$ is the \acrshort{td} error $\delta_t$, measuring the difference between the current estimate $Q\left(S_t, A_t\right)$ and the bootstrapped target $R_{t+1} + \gamma \max_a Q\left(S_{t+1}, a\right)$. The target is said to be `bootstrapped' because it relies on the current estimate $Q\left(S_{t+1}, a\right)$ rather than a complete return computed from all future rewards to the episode's end. The $\max$ operator makes the update off-policy because the bootstrap target always uses the greedy action in $S_{t+1}$, irrespective of the action actually selected by the behavior policy \citep[pp.~131, 257]{suttonReinforcementLearningIntroduction2020}. Being a \acrshort{td} algorithm, Q-learning learns from sample data and does not require a model of the environment's dynamics \citep[p.~119]{suttonReinforcementLearningIntroduction2020}. When the state and action spaces are finite and small enough for every entry of $Q$ to be visited repeatedly, tabular Q-learning converges to $q_*$ under standard stochastic approximation conditions \citep[p.~33]{suttonReinforcementLearningIntroduction2020}.

These convergence conditions require that every state-action pair be visited infinitely often and that the step size decreases appropriately \citep[p.~131]{suttonReinforcementLearningIntroduction2020}. Under function approximation, individual state-action pairs no longer have independent storage. Instead, a shared parameter vector $\theta$ determines the estimate for all states simultaneously, and updating one estimate necessarily affects others. This coupling means that the convergence proof for the tabular case does not carry over.

Pong has a state represented by 84x84x4 stack of pixel intensities which is enormously larger than any tractable table \citep[p.~534]{mnihHumanlevelControlDeep2015}. To learn action values directly from such high-dimensional input \citet[p.~529]{mnihHumanlevelControlDeep2015} developed a deep Q-network that uses a deep convolutional network \citep{lecunGradientbasedLearningApplied1998}. Convolutional networks exploit the spatial structure inherent in image data through local receptive fields and shared weights, making them a natural choice for processing raw pixel observations. The learned action-value function is parameterized as $Q\left(s,a;\theta_i\right)$ where $\theta_i$ are the weights of the Q-network at iteration $i$. The bootstrapped target is represented as $r + \gamma \max_{a^{\prime}} Q\left(s^{\prime}, a^{\prime}; \theta_i\right)$. When computing the gradient of the loss, the target is treated as a fixed quantity even though it depends on $\theta_i$, so the gradient is taken only through the estimate $Q\left(s,a;\theta_i\right)$. This is known as a semi-gradient method \citep[p.~202]{suttonReinforcementLearningIntroduction2020}. The deep Q-network is trained by minimizing the following loss function via \acrfull{sgd} at each iteration $i$:
\begin{equation}\label{eq:q_learning_loss_function}
	L_i\left(\theta_i\right) = \mathbb{E}_{s, a, r, s^{\prime}} \Bigg[\bigg(r + \gamma \max_{a^{\prime}} Q(s^{\prime}, a^{\prime};\theta_i) - Q(s,a;\theta_i)\bigg)^2\Bigg].
\end{equation}
This formulation, however, is unstable in practice. The target depends on the same parameters being updated, and the training data is generated sequentially by the agent's own interaction with the environment. Section~\ref{sec:stabilization} addresses both problems.

\section{Stabilization mechanisms}
\label{sec:stabilization}

The combination of function approximation, bootstrapping and off-policy training is known to risk instability and divergence in what \citet[p.~264]{suttonReinforcementLearningIntroduction2020} calls the deadly triad. \acrshort{dqn} exhibits all three where a neural network approximates $Q$, the target bootstraps from the network's own estimate, and the max operator evaluates the greedy policy regardless of which action generated the transition. In practice, this instability manifests in two concrete ways. 

First, a correlation exists between the sequence of observations \citep[p.~529]{mnihHumanlevelControlDeep2015}, since input frames of Pong within a short time-frame are quite similar. \acrshort{sgd} on correlated input samples can oscillate or converge to poor solutions because each mini-batch of input data is locally biased rather than representative of the full state space. Experience replay addresses this by storing transitions and sampling them uniformly at random, breaking the temporal structure. Beyond decorrelation, experience replay also improves sample efficiency. Each transition can contribute to multiple weight updates rather than being used once and discarded, allowing the network to extract more learning signal from a fixed amount of environment interaction \citep[p.~535]{mnihHumanlevelControlDeep2015}.

Second, a dependency exists between the action-values $Q(s,a;\theta_i)$ and the target values $r+\gamma \max_{a^{\prime}} Q\left(s^{\prime}, a^{\prime};\theta_i\right)$ because they share the same parameters $\theta_i$ in \eqref{eq:q_learning_loss_function}. This creates a moving target that shifts with each gradient step. \acrshort{dqn} addresses this by providing a fixed set of parameters $\theta_i^{\,\text{--}}$ for computing that target, updated only every C steps using a hard copy of $\theta_i$ parameters of the Q-network. This converts the continuously moving regression target into a piecewise-stationary one, giving the network a stable objective to optimize between successive copies.
	
The full \acrshort{dqn} loss function in \citet[p.~529]{mnihHumanlevelControlDeep2015} is:
\begin{equation}\label{eq:q_learning_loss_function_in_dqn}
	L_i\left(\theta_i\right) = \mathbb{E}_{(s, a, r, s^{\prime}) \sim U(D)} \Bigg[\bigg(r + \gamma \max_{a^{\prime}} Q(s^{\prime}, a^{\prime};\theta_i^{\,\text{--}}) - Q(s,a;\theta_i)\bigg)^2\Bigg],
\end{equation}
where $(s, a, r, s^{\prime}) \sim U(D)$ denotes that the mini-batch samples are uniformly drawn from the agent's stored experience dataset $D$.

The policy learned by \acrshort{dqn} is derived implicitly from the learned Q-values via greedy action selection, where the agent always takes the action with the maximum estimated value. This policy is deterministic, and exploration is imposed externally through an $\epsilon$-greedy mechanism that selects a uniformly random action with probability $\epsilon$. While effective, this treats all non-greedy actions as equally worth exploring and ignores any learned information about their relative value. Chapter~\ref{chap:a2c} develops the alternative approach, in which the policy is parameterized directly as a probability distribution over actions and optimized via gradient ascent.